\section{Conditional expectation and conditional variance}
\label{sec:conditional}
\noindent\textbf{Reading.} Freund, Chapter~4, \S8, pp.~138--139, Definition~10 and Example~20. The final subsection on averaging conditional means and variances is supplementary reading from Bertsekas--Tsitsiklis, \S4.3, chapter-local pp.~17--24; it strengthens the listed exam topics without replacing the direct calculations required by Homework~5.

\subsection{Definition: take an expectation under the conditional law}
Once $Y=y$ is given, use the distribution of $X\mid Y=y$ as the probability weights. For a function $g$,
\[
\boxed{\E[g(X)\mid Y=y]=\sum_x g(x)p_{X\mid Y}(x\mid y)}
\]
in the discrete case, or
\[
\boxed{\E[g(X)\mid Y=y]=\int g(x)f_{X\mid Y}(x\mid y)\,\dd x}
\]
in the continuous case. The limits are the conditional support. Treat $y$ as a fixed parameter while summing or integrating over $x$. The choices $g(x)=x$ and $g(x)=x^2$ give the conditional first and second raw moments.

Substituting the conditional density formula gives a convenient quotient:
\[
 \E[g(X)\mid Y=y]
 =\frac{\displaystyle\int g(x)f_{X,Y}(x,y)\,\dd x}{f_Y(y)},
 \qquad f_Y(y)>0.
\]
Use a sum and $p_Y(y)>0$ in the discrete version. This formula can save writing a separate conditional density, but the denominator and support must still be correct. The numerator alone is generally not a conditional expectation.

\subsubsection*{Worked example: a conditional mean from a table}
Suppose
\[
\begin{array}{c|cc|c}
 &Y=0&Y=1&p_X(x)\\\hline
X=1&1/8&3/8&1/2\\
X=2&3/8&1/8&1/2\\\hline
p_Y(y)&1/2&1/2&1
\end{array}
\]
Given $Y=0$, divide the first column by $1/2$, obtaining conditional probabilities $1/4,3/4$. Thus
\[
 \E[X\mid Y=0]=1\left(\frac14\right)+2\left(\frac34\right)=\frac74,
 \qquad
 \E[X^2\mid Y=0]=\frac14+4\left(\frac34\right)=\frac{13}{4}.
\]
Given $Y=1$, the conditional probabilities reverse, so $\E[X\mid Y=1]=5/4$.

\paragraph{Checkpoint.} If $P(X=0\mid Y=2)=1/3$ and $P(X=3\mid Y=2)=2/3$, find $\E[X\mid Y=2]$ and $\E[X^2\mid Y=2]$. \emph{Answer:} $2$ and $6$. Squaring the conditional mean would give $4$, which is a different quantity.

\subsection{Conditional variance and support-dependent calculations}
The conditional mean is $m(y)=\E[X\mid Y=y]$. The conditional variance measures spread around that conditional mean:
\[
\boxed{\Var(X\mid Y=y)
 =\E[X^2\mid Y=y]-\{\E[X\mid Y=y]\}^2.}
\]
This follows by expanding $\E[(X-m(y))^2\mid Y=y]$, treating $m(y)$ as constant. It is not $\E[X^2\mid Y=y]-\{\E[X]\}^2$: use the same conditional law for both terms.

For the preceding table,
\[
 \Var(X\mid Y=0)=\frac{13}{4}-\left(\frac74\right)^2=\frac3{16}.
\]
Given $Y=1$, the variance is also $3/16$, although the mean differs. Equal conditional variances do not imply independence.

\subsubsection*{Worked example: a conditional density on a moving interval}
Let $f(x,y)=3y$ on $0<x<y<1$, zero elsewhere. Integrate over the slice in $x$:
\[
 f_Y(y)=\int_0^y3y\,\dd x=3y^2,
 \qquad f_{X\mid Y}(x\mid y)=\frac1y,\quad0<x<y,
 \quad0<y<1.
\]
Thus $X\mid Y=y$ is uniform on $(0,y)$. Direct calculation gives
\[
 \E[X\mid Y=y]=\int_0^y\frac{x}{y}\,\dd x=\frac y2,
 \qquad
 \E[X^2\mid Y=y]=\int_0^y\frac{x^2}{y}\,\dd x=\frac{y^2}{3},
\]
and therefore
\[
 \Var(X\mid Y=y)=\frac{y^2}{3}-\frac{y^2}{4}=\frac{y^2}{12}.
\]
At $y=1/2$, the conditional mean is $1/4$ and the conditional variance is $1/48$. Replacing the upper limit $y$ by $1$ would fail to normalize the conditional law and produce incorrect moments.

\paragraph{Checkpoint.} In this triangular model, find $\E[X^2\mid Y=y]$ and $\Var(X^2\mid Y=y)$. \emph{Answer:} $y^2/3$ and $y^4/5-(y^2/3)^2=4y^4/45$; the second calculation needs the conditional fourth moment of $X$.

\subsection{A function of a value versus a random variable}
The expression $\E[X\mid Y=y]=m(y)$ is a number once $y$ is fixed. The expression $\E[X\mid Y]=m(Y)$ is a random variable, because its value depends on the unobserved outcome of $Y$.

In the triangular example, $m(y)=y/2$, so $\E[X\mid Y]=Y/2$, while $\E[X\mid Y=1/2]=1/4$. In the table example, $m(Y)$ takes values $7/4$ and $5/4$, each with probability $1/2$.

Conditional expectation remains linear:
\[
 \E[aX+b\mid Y=y]=a\E[X\mid Y=y]+b.
\]
Likewise $\Var(aX+b\mid Y=y)=a^2\Var(X\mid Y=y)$. A quantity determined by the given variable can be treated as fixed inside the conditional expectation; for example,
$\E[Yg(X)\mid Y=y]=y\E[g(X)\mid Y=y]$. With independence and integrability, $\E[X\mid Y]=\E[X]$ almost surely. A constant conditional mean alone does not prove independence, since higher conditional moments or other distributional features may still change.

\paragraph{Checkpoint.} If $\E[X\mid Y=y]=2y-1$, find $\E[3X+4\mid Y=2]$ and express $\E[X\mid Y]$. \emph{Answer:} $13$ and $2Y-1$.

\subsection{Supplement: averaging conditional means and variances}
The law of total expectation, also called the tower rule, says
\[
 \boxed{\E[X]=\E[\E[X\mid Y]].}
\]
In a continuous model, this is $\E[X]=\int m(y)f_Y(y)\,\dd y$. The reason is that the conditional expectation's numerator, multiplied by $f_Y(y)$ and integrated in $y$, reconstructs the joint-law integral for $\E[X]$.

For finite second moments, the law of total variance is
\[
\boxed{\Var(X)=\E[\Var(X\mid Y)]+\Var(\E[X\mid Y]).}
\]
The first term is the average spread within conditional groups; the second is the spread between their means. To derive it, average the conditional identity
$\Var(X\mid Y)=\E[X^2\mid Y]-m(Y)^2$ and add
$\Var(m(Y))=\E[m(Y)^2]-\{\E[X]\}^2$.

\subsubsection*{Worked check: recover unconditional moments}
For $f(x,y)=3y$ on $0<x<y<1$, $f_Y(y)=3y^2$, so
\[
 \E[Y]=\frac34,\qquad\E[Y^2]=\frac35,\qquad\Var(Y)=\frac3{80}.
\]
Using the conditional results,
\[
 \E[X]=\E[Y/2]=\frac38,
\]
\[
 \Var(X)=\E[Y^2/12]+\Var(Y/2)
 =\frac1{20}+\frac3{320}=\frac{19}{320}.
\]
The direct check is $\E[X^2]=\E[Y^2/3]=1/5$, and $1/5-(3/8)^2=19/320$.

\paragraph{Checkpoint.} In the discrete table above, recover $\E[X]$ and $\Var(X)$ using the two total laws. \emph{Answer:} $\E[X]=(7/4+5/4)/2=3/2$. The average conditional variance is $3/16$ and the variance of the two conditional means is $1/16$, so $\Var(X)=1/4$.

\paragraph{Textbook practice.} Freund Chapter~4, Example~20 (p.~139), gives the direct conditional-mean and conditional-variance workflow used in Homework~5, Problem~7. Exercise~60 (p.~140) conditions a univariate law on an interval; it is useful additional practice in recognizing the conditioning event and normalizing correctly.

